% SYNTHETIC TEST UNIT (integration tools) -- deliberate errors are marked "DELIBERATE"
\chapter{引言（合成测试）}\label{ch:intro}
本章是集成工具测试用的合成文本。文献待补：某某 2020。% DELIBERATE K9 printed planning marker
本书不引用参考文档 factorization\_qcd-2.pdf。% DELIBERATE K11 mention of the reference document

\section{严格性标签}\label{sec:intro-tags}

\begin{definition}[严格性标签]\label{def:intro-tags}
五个标签 \tagR、\tagA、\tagP、\tagC、\tagV（合成测试）。
% uses: none
\end{definition}

\begin{assumption}[QCD 的 Wightman 结构 \tagC]\label{asm:intro-wightman}
(W1)--(W4)（合成测试）。
地位：公开问题。
文献：无。
被使用于：\cref{thm:dis-moments}。
解除途径：构造 QCD。
\uses{def:intro-tags}
\end{assumption}

\begin{assumption}[非微扰算符乘积展开 \tagC]\label{asm:intro-np-ope}
(O1)(O2)（合成测试）。
地位：公开问题。
文献：无。
被使用于：\cref{thm:dis-moments}。
解除途径：非微扰定义。
% uses: none
\end{assumption}

\begin{assumption}[桥接假设 \tagC]\label{asm:intro-bridge}
(Br1)--(Br3)。
地位：公开问题。
文献：无。
被使用于：\cref{thm:dy-factorization,prop:dy-npline}。
解除途径：无。
% uses: none
\end{assumption}

\begin{proposition}[代数命题 \tagR]\label{prop:intro-algebra}
（$R_{\mathrm{def}}$：由算符定义直接推出。）合成测试。（本命题只涉及部分子态。）
\uses{def:intro-tags}
\end{proposition}
